\section[What are BSSs?]{What are \acp{bss}?}
\label{sec:what_are_bss}

\acp{bss} are mobility services that provide access to shared bicycles for short periods of time, enabling users to complete individual trips without owning a bicycle~\cite{Fishman2016}. The principle is simple: pick up a bike, ride it to the desired destination, and leave it for the next user. By removing the burdens of ownership—such as storage, maintenance, and theft risk—\acp{bss} make cycling more accessible to a wider range of users. This is particularly relevant given the persistent gap between bicycle availability and actual usage: despite relatively high ownership rates in many regions, cycling still accounts for less than 5\% of daily trips in most countries~\cite{Chen2022}. By offering convenient, low-commitment access, \acp{bss} aim to convert this latent cycling potential into actual behaviour and position cycling as a practical alternative to private motorized travel. At the same time, these systems offer cities a scalable and visible tool to advance broader goals in sustainability, public health, and urban livability~\cite{Shaheen2010, Ricci2015}.

\acp{bss} can be implemented through different operational and technological models, but most contemporary systems share four core components: (i) a fleet of bicycles; (ii) pick-up and return infrastructure, such as docking stations or designated parking areas; (iii) user access and payment mechanisms; and (iv) operational systems for maintenance, monitoring, and redistribution.

Building on this common structure, \acp{bss} vary widely in how they operate and are configured. A fundamental distinction is between \textbf{docked} (station-based) and \textbf{dockless} (free-floating) systems. Docked systems require users to pick up and return bicycles at designated docking stations, whereas dockless systems allow bikes to be parked freely within defined areas using built-in locks and GPS technology. Beyond this typology, cities implement a range of additional design choices, including whether fleets consist of \acp{m-b}, \acp{e-b}, or both; how pricing schemes combine pay-per-use and subscription options; and whether incentives are introduced to encourage short trips or fleet turnover. Despite this diversity, nowadays most \acp{bss} rely on real-time monitoring, preventive maintenance, and active fleet rebalancing to ensure reliable service.

In addition to their technical and operational characteristics, \acp{bss} also differ in how they are owned, managed, and funded, with ownership and management models typically grouped into three broad modalities~\cite{Li2021, Cao2022}. Public systems are fully owned and managed by public organisations. Madrid’s \textit{BiciMAD}, run directly by the city’s municipal transport company EMT, is a clear example~\cite{BiciMAD2025}. At the other extreme, private systems are owned and operated by companies such as \textit{Lime} and \textit{Donkey Republic}. Between these two poles, hybrid public–private partnerships are common, with ownership remaining in public hands while operations are outsourced to private firms. Barcelona’s \textit{Bicing}, owned by the municipality but operated under concession by \textit{Pedalem Barcelona}~\cite{BicingPedalem2017}, illustrates this model. Funding structures reflect these arrangements: public systems are typically supported by municipal budgets or subsidies, private systems rely on user fees and advertising revenues, and hybrid models combine both sources~\cite{Shaheen2010}.

% Cortez-Ordoñez2024 Are we back to normal?-> It was initially operated by Clear Channel until 2018; since 2019, Pedalem Barcelona has been the company responsible for providing the service under the supervision of the public company Barcelona de Serveis Municipals (BSM). 
%  Bike Sharing Systems in Munich: A (non)users’ behavioural, socio–demographic and psychographic analysis. Apartat 2.1.1.4 Sobre ownership
% RAKY_JULIO_CASTILLO_Explanatory factors for improving performance of bike sharing systems Secció 2.4. Business models for BSS
% Midgley2011 -> Taula 8: A comparison of operating and financing structures

In the following section, the historical evolution of \acp{bss} is presented, tracing their development from early experiments in the 1960s to today’s digitally managed fleets of connected bicycles operating in cities around the world.